\section{Matching without a supplied root mark}
\label{sec:v78-root}

The local prefix argument proves more than local solvability once one
uses determinism at every zero of the matching equation.  The following
one-dimensional fact is useful.

\begin{lemma}[Downward zeros]\label{lem:v78-downward}
Let $h\in C^1([a,b])$.  If $h'(x)<0$ at every zero in $(a,b)$,
then there is at most one such zero.  If the endpoint values are nonzero,
a zero exists exactly when $h(a)>0>h(b)$.
\end{lemma}
\begin{proof}
A zero is isolated, positive immediately to its left and negative
immediately to its right.  If there were two zeros, take a point just
after the first with negative value and a point just before the second
with positive value.  Between them the zero set is compact and discrete,
and hence finite.  Every crossing there is from positive to negative,
which cannot carry a negative starting value to a positive ending value.
This is a contradiction.  The assertion about endpoints follows from
the intermediate value theorem and the signs on either side of the
unique zero.
\end{proof}

\begin{theorem}[Root-selection-free prefix reconstruction]
\label{thm:v78-root-free}
Let $U(s,t)$ and $V(s,v)$ be regular actions for a physical reflection
word and its one-flight extension.  Suppose that both are defined on
products with a common connected initial interval $[a,b]$, including
its endpoints.  For fixed $(t,v)$ put
\[
 H(s,t,v)=V_s(s,v)-U_s(s,t).
\]
Every interior zero corresponds to the same physical prefix and obeys
\begin{equation}\label{eq:v78-root-sign}
 H_s=-\frac{U_{st}^2}{B}<0,\qquad B=U_{tt}+\ell_{tt}>0.
\end{equation}
There is at most one interior root.  On the data-determined open set
\[
 \Omega_{a,b}=\{(t,v):H(a,t,v)>0>H(b,t,v)\}
\]
it exists, is smooth, and determines
\[
 \ell(t,v)=V(\sigma(t,v),v)-U(\sigma(t,v),t).
\]
No root label, small root bracket or measured initial momentum is needed.
Every regular extended trajectory with initial coordinate in $(a,b)$
and terminal coordinates in the interiors of the two gates belongs
locally to such an open set after restriction within the original gates.
\end{theorem}
\begin{proof}
At a zero, the two trajectories start at the same boundary point and
have the same initial tangential covector.  The prescribed outgoing side
and unit speed determine the outgoing velocity uniquely.  Determinism
of nongrazing billiard motion therefore identifies the entire prefix.
In particular, every zero is physical, not just the zero chosen near a
reference trajectory.  The stationary composition identity and positive
interior Hessian from Lemma~\ref{lem:v77-regular-action} apply there.
Eliminating prefix variables gives the positive Schur complement $B$;
differentiation gives~\eqref{eq:v78-root-sign}.

Apply Lemma~\ref{lem:v78-downward} to $s\mapsto H(s,t,v)$.  It is
unnecessary, and in general unjustified, to assert $H_s<0$ at points
which are not roots.  The endpoint-sign criterion, the implicit function
theorem at the unique root, and the stationary composition identity give
all the stated conclusions.  For a regular matched trajectory, take
$a,b$ strictly on its two sides in the original common gate.  The signs
hold there and persist for nearby terminal coordinates.  Rational
choices are possible by openness.  These are constructed from recovered
actions; they are not extra observations.
\end{proof}

\begin{corollary}[Data-derived finite certificates]
\label{cor:v78-certificates}
In the finite-atlas theorem, supplied matching brackets and supplied
clock-anchor points can be omitted.  For an exact admissible law tuple,
usable clock anchors, root domains, distance anchors and overlap triples
can be selected by strict tests on the recovered functions and the fixed
scalar labels.  A finite collection suffices whenever an identifying
atlas exists.
\end{corollary}
\begin{proof}
Nonconstancy and Theorem~\ref{thm:v77-clock} give a nonzero anchor
denominator on an open set, so a countable dense list of pairs of scalar
points contains a usable pair.  Recover the actions.  Enumerate rational
common initial subintervals and rational terminal rectangles with closures
inside the open sets $\Omega_{a,b}$ of the preceding theorem.  Each true
local matching patch has such a restriction.  On it recover $\ell$.
The nonzero finite determinants in
Proposition~\ref{prop:v77-distance} are open evaluation tests, as are
noncollinearity of overlap triples.  Rational parameters therefore
suffice.  The scalar cover and its graph are combinatorial tests on these
intervals.  The finite cover from Theorem~\ref{thm:v77-finite-atlas}
ensures that one finite selection passes all tests.

This is a law-level selection statement.  With computable $C^r$ names
and certified derivative bounds, each strict inequality on a compact
rectangle can be semidecided by a rational grid and its derivative error
bound; dovetailing these tests then terminates on the stated class.
Without a representation of functions, exact laws must not be called
a finite computational input.  No test here certifies physical
admissibility of an arbitrary clock request outside the assumed class.
\end{proof}

\subsection{A principle for scalar generating systems}

The argument is not specific to the Euclidean chord.  Suppose a scalar
deterministic twist system has smooth generating actions $U,V$ for a
word and its extension, with nonzero $U_{st}$ and $V_{sv}$.  Suppose equality of their
initial covectors identifies the common trajectory, and the stationary
composition
\[
             V(s,v)=\operatorname*{stat}_t\{U(s,t)+L(t,v)\}
\]
has $B=U_{tt}+L_{tt}>0$ at every match.  Supply three normalized
clocked laws for each action with an unknown positive common weight.

\begin{theorem}[Clocked generating-action reconstruction]
\label{thm:v78-generating}
Under these hypotheses, the laws determine the absolute one-step
generating function $L$ on every matched open domain.  They also determine
the corresponding local canonical transformation, without a momentum
observation or a root-selection mark.  On bounded nested domains with
positive clock and matching margins, recovery of $L$ is locally
Lipschitz from $C^{k+2}$ laws to $C^k$ functions for every finite $k$.
\end{theorem}
\begin{proof}
The clock theorem recovers $U,V$ and their constants.  At any zero of
$V_s-U_s$, the deterministic prefix property makes the stationary
composition applicable.  Differentiating it gives
$V_{ss}-U_{ss}=-U_{st}^2/B$.  Lemma~\ref{lem:v78-downward} proves
uniqueness on a connected common interval.  Subtraction at the root
recovers $L$, including its constant.  Differentiation yields its
canonical relation.  The implicit differentiation and composition
estimate in Proposition~\ref{prop:v77-cancellation} uses only these
margins, so proves the stated finite-order bound without invoking
Euclidean geometry.
\end{proof}
